\par\Needspace{9\baselineskip}\bigskip
\pdfbookmark[3]{Darstellung, 19}{pallas-II-article19}
\begin{center}19.\end{center}
\begingroup
\setlength{\lineskip}{4pt}\setlength{\lineskiplimit}{4pt}
\noindent\hspace*{\parindent}Wendet man diese Integrationsmethoden auf die Änderungen unserer Elemente an, so erhält man, abgesehen von den bei der Integration hinzukommenden Konstanten, die Werte von \(i,\ \AscendingNode,\ \varphi,\ n,\ \tilde\omega,\ \epsilon\) zu den Epochen, die zwischen jenen liegen, für die man die momentanen Änderungen berechnet hat, und den Wert von \(\int n\mathrm{d}t\) zu eben diesen ursprünglichen Epochen, wenn man die Formeln I, II verwendet, oder umgekehrt, wenn man die Formeln III, IV vorzieht. Man kann auch für die einfache Integration Formel III und für die doppelte Integration Formel II verwenden. Damit erhält man alle Elemente zu denselben Epochen. Im Übrigen ist die Wahl ziemlich gleichgültig, da diese Werte der Elemente nur dazu dienen, mittels der bekannten Interpolationsmethoden die Werte zu ermitteln, die sich auf die Epochen der verschiedenen Beobachtungen beziehen. Die Auswertung dieser Beobachtungen liefert nach bekannten Methoden die Werte der Konstanten. Es sind ihrer sieben, die sich aber auf sechs reduzieren, da die Konstante \(\alpha\), die sich aus \(\int\dfrac{\mathrm{d}\epsilon}{\mathrm{d}t}\mathrm{d}t\) ergibt, und die Konstante \(\beta\), die aus \(\int n\mathrm{d}t\) oder \(\int\!\int\dfrac{\mathrm{d}n}{\mathrm{d}t}\mathrm{d}t^2\) hervorgeht, zu einer einzigen zusammenfallen.
\par\endgroup
